\runninghead{Each Text in Its Setting}
\section{Each Text in Its Setting}
\label{sec:each-element}

\subsection{Introit: the soul lifted to God}
\label{sec:element-introit}

Psalm 24 (25) is a prayer for guidance and forgiveness, headed in the Douay
``Unto the end, a psalm for David.'' After the three verses the Introit sings,
the psalmist asks to be taught, confesses ``the sins of my youth and my
ignorances,'' and begs God to remember ``thy bowels of compassion; and thy
mercies that are from the beginning of the world'' (vv.~5--7). He says that
``all the ways of the Lord are mercy and truth'' (v.~10), that his eyes are
``ever towards the Lord: for he shall pluck my feet out of the snare''
(v.~15), and he ends by praying for the whole people: ``Deliver Israel, O God,
from all his tribulations'' (v.~22). The man who lifts his soul in the first
verse is in danger, burdened by his sins and mocked by enemies.

The chant keeps three of his acts. He lifts his soul, he trusts, and he
appeals to a rule he has learned: none who wait for God are put to shame.
St Augustine, whose own psalter read \latin{qui te
exspectant} as the chant does, hears in the psalm ``Christ \dots\ in the
person of the Church,'' and gives the lifting a sense: ``with spiritual
longing have I lift up the soul, that was trodden down on the earth with
carnal longings'' (\work{Enarrationes in Psalmos} 24.1--2). The enemies who
would laugh are, for him, the ``serpent-like and secret suggestions'' that
flatter a man to his fall with ``Well done, well done'' (24.3).

The verse asks God to show his ways and teach his paths, the request the
Gradual will repeat in other words.

\subsection{Collect: Come, and rescue us from our sins}
\label{sec:element-collect}

The Collect is short and urgent. It asks the Lord to rouse his power,
\latin{Excita \dots\ poténtiam tuam}, and then simply \latin{et veni}, ``and
come.'' Its conclusion, \latin{Qui vivis et regnas cum Deo Patre}, shows that
it is addressed to the Son. The purpose clause names what is feared:
\latin{ab imminéntibus peccatórum nostrórum perículis}, ``from the
threatening dangers of our sins.'' The dangers hang over us
(\latin{imminéntibus}), and they are ours, made by our sins. Two pairs of
verbs follow, each with an ablative of Christ's action: \latin{te protegénte
éripi}, to be snatched away with him protecting, and \latin{te liberánte
salvári}, to be saved with him setting free. Schuster paraphrases the petition
as a request that our Lord ``save us by his power from the dangers to which
we are exposed through our sins'' (\work{The Sacramentary}, vol.~1, p.~321).

\subsection{Epistle: the hour to rise}
\label{sec:element-epistle}

The lesson comes from the moral part of St Paul's letter to the Romans. Just
before it he has told them to ``owe no man any thing, but to love one
another,'' since ``love therefore is the fulfilling of the law'' (Rom
13:8, 10). The Vulgate then continues \latin{Et hoc scientes tempus}, ``And
that, knowing the season'': the love just commanded is to be practised by
people who know what time it is. The Missal's opening, \latin{Fratres:
Sciéntes}, drops the link and lets the lesson stand alone.

Paul's argument has three steps. It is the hour to rise from sleep, because
salvation is nearer than when we first believed. The night is past and the
day at hand, so the works of darkness are to be thrown off and the armour of
light put on. To walk becomingly as in the day means to leave six things,
grouped in three pairs: revelling and drunkenness, the bed and wantonness,
strife and envy. The last command is \latin{induímini Dóminum Iesum
Christum}, ``put ye on the Lord Jesus Christ.''
The verbs of clothing run through the passage: \latin{abiciámus},
\latin{induámur}, \latin{induímini}. The Missal stops before the second half
of verse~14, so the lesson ends on the name of Christ.

St Augustine read these verses in a garden and was converted. He
tells the story in the \work{Confessions}: a child's voice from a neighbouring
house chanting ``Take up and read; take up and read,'' and his opening the
book of the Apostle at ``Not in rioting and drunkenness \dots\ but put ye on
the Lord Jesus Christ, and make not provision for the flesh, to fulfil the
lusts thereof.'' ``No further would I read, nor did I need''
(\work{Confessions} VIII.12.29). Schuster says that ``the Church rouses us at
last from sleep by those inspired words of the Apostle which once wrought
the conversion of Augustine'' (p.~321).

\subsection{Gradual: those who wait, and the ways made known}
\label{sec:element-gradual}

The Gradual returns to Psalm 24 after the Epistle. Its respond is the end of
verse~3, the rule the Introit cited: \latin{Univérsi, qui te exspéctant, non
confundéntur, Dómine}, ``None who wait for thee shall be confounded, O
Lord.'' The chant adds the final \latin{Dómine}, so that the rule becomes an
address. Bellarmine explains why the psalmist can state it: ``we have learned
by long experience, from the examples of our ancestors, and from your own
promises, that those who put their trust in you, and patiently expect your
help, were never disappointed''; to wait on the Lord ``means to expect him in
the certain hope of assistance'' (\work{Explanatio}, Ps 24, on vv.~2--3).

The verse repeats the Introit's request with a different verb: \latin{notas
fac mihi}, ``make known to me,'' where the Introit has \latin{demónstra}.
Augustine's psalter had the same words as the Gradual, which the English of
his \work{Enarrationes} renders ``Make Thy ways, O Lord, known to me''; he
calls the paths ``narrow, and known to few'' (24.4). After the Epistle's
list of ways to walk and not to walk, the Gradual asks God to make his own way
known.

\subsection{Alleluia: Shew us thy mercy}
\label{sec:element-alleluia}

Psalm 84 (85) begins with thanks for a restoration already given: ``Lord, thou
hast blessed thy land: thou hast turned away the captivity of Jacob. Thou hast
forgiven the iniquity of thy people'' (vv.~2--3). Then the psalmist asks for
it again, as if the anger had returned: ``Convert us, O God our saviour: and
turn off thy anger from us. Wilt thou be angry with us for ever?''
(vv.~5--6). The Alleluia verse is the climax of that prayer: ``Shew us, O
Lord, thy mercy; and grant us thy salvation'' (v.~8). In the next verses the
psalmist listens for God's answer and hears a promise of peace: ``Surely his
salvation is near to them that fear him: that glory may dwell in our land.
Mercy and truth have met each other: justice and peace have kissed''
(vv.~10--11).

The word for salvation in v.~8 is \latin{salutáre}, a neuter adjective used as
a noun, ``thy saving thing.'' Augustine reads it as a person: \latin{salutare
tuum, Christum tuum}, ``thy salvation, thy Christ'' (\work{En.\ in Ps.}
84.9). He reasons from the psalm that Christ, already given, is still to be
asked for.

\subsection{Gospel: the signs and the coming of the Son of man}
\label{sec:element-gospel}

The lesson is the end of Christ's last discourse in St Luke, spoken in the
Temple in the days before his Passion (Lk 21:37--38). It began when some
admired the Temple's stones and he answered that ``there shall not be left a
stone upon a stone'' (vv.~5--6). The disciples asked when, and by what sign
(v.~7). He warned them of false christs, of wars and earthquakes and ``terrors
from heaven,'' of persecution in which ``a hair of your head shall not
perish'' (vv.~8--18), and then of Jerusalem compassed with armies, its people
led captive and the city ``trodden down by the Gentiles till the times of the
nations be fulfilled'' (vv.~20--24). The Vulgate's next verse begins
\latin{Et erunt signa}; the Missal begins the lesson at the signs with its own
formula, \latin{In illo témpore}. Augustine and St Cyril of Alexandria both
take Luke to pass at this point from the fall of the city to the end:
Augustine observes that Luke keeps the siege of Jerusalem apart from the last
coming (Letter 199.27), and Cyril writes that Christ ``transfers His words
from this subject to the time of the consummation'' (\work{Commentary on
Luke}, sermon 139).

The lesson has three parts. First come the signs in the heavens, the distress
of nations, men withering for fear, and the powers of heaven moved; then the
coming of the Son of man \latin{in nube}, ``in a cloud,'' with great power and
majesty; then the command to look up and lift the head, because redemption is
near. A parable follows: when the fig tree and all the trees put out their
fruit, summer is near, and when these things happen the kingdom of God is
near. The lesson ends with two sayings sealed by \latin{Amen}: this
generation will not pass till all is done, and heaven and earth will pass but
Christ's words will not. Just after the lesson Christ warns against hearts
``overcharged with surfeiting and drunkenness and the cares of this life''
(v.~34), and commands, ``Watch ye, therefore, praying at all times''
(v.~36).

St Gregory the Great preached on this Gospel in Rome. His text read
\latin{in nubibus}, ``in the clouds,'' where the Missal has the singular. In
the course of his homily he quotes the prophet Sophonias on the day of the
Lord: \latin{Dies irae, dies illa}, ``That day is a day of wrath'' (Zeph
1:15, \work{Homiliae in Evangelia} I.1.6). Schuster records that the later
sequence that begins with the same words was once sung on this Sunday before
the Gospel (p.~323).

\subsection{Offertory: the Introit sung again}
\label{sec:element-offertory}

At the offering the choir sings the Introit's antiphon again, word for word,
without the psalm verse. Schuster calls Psalm 24 ``the typical psalm of the
first Sunday of Advent'' and finds the reason for its choice in the rule of
verse~3, which he paraphrases: ``He who with watchful trust waits upon the
Lord shall never be confounded'' (p.~322). Psalm 24 is sung three times in
this Mass: at the entrance, between the lessons and at the offering.

\subsection{Secret: purer to their source}
\label{sec:element-secret}

The Secret's subject is \latin{Hæc sacra}, ``these holy things,'' the gifts on
the altar. It asks that they make us, \latin{poténti virtúte mundátos},
``cleansed by their mighty power,'' come \latin{purióres}, ``purer,''
\latin{ad suum princípium}, to their own beginning or source. The gifts
cleanse, and they lead those they cleanse back to the one from whom they
come. Schuster finds here ``the ancient idea that runs through the Roman
epiklesis'': the sacrifice's virtue is asked to purify our minds ``that we may
worthily receive him who is its Author'' (p.~322).

\subsection{Communion: the Lord will give goodness}
\label{sec:element-communion}

The Communion takes the verse that closes the promise of Psalm 84: ``The
Lord will give goodness: and our earth shall yield her fruit'' (v.~13). In the
psalm it follows ``Truth is sprung out of the earth: and justice hath looked
down from heaven'' (v.~12), and it is followed by ``Justice shall walk before
him'' (v.~14). The psalm moves between heaven's gift and the earth's yield,
and the Communion keeps both halves. The Missal omits the Vulgate's opening
\latin{Etenim}, ``For,'' so that the chant states the promise without its
link to the verse before. Augustine read \latin{suavitatem}, ``sweetness,''
where the Missal and the Vulgate have \latin{benignitátem}, ``goodness''
(\work{En.\ in Ps.} 84.15). Who or what the fruitful earth is divides the
commentators.

\subsection{Postcommunion: mercy received in the temple}
\label{sec:element-postcommunion}

The Postcommunion borrows the words of Ps 47:10 (48:10), \latin{Suscepimus,
Deus, misericordiam tuam in medio templi tui}, ``We have received thy mercy,
O God, in the midst of thy temple,'' as Schuster notes (p.~323), and turns
them into a request, \latin{Suscipiámus}, ``may we receive.'' The purpose
clause looks forward: \latin{ut reparatiónis nostræ ventúra solémnia cóngruis
honóribus præcedámus}, ``that we may go before the coming solemnities of our
restoration with fitting honours.'' Schuster takes the solemnities to come
to be Christmas, ``the approaching feast of our common ransom'' (p.~323).
